\documentclass[11pt,a4paper]{article}
\usepackage[T1]{fontenc}
\usepackage[utf8]{inputenc}
\usepackage[english]{babel}
\usepackage{lmodern}
\usepackage{microtype}
\usepackage[margin=2.3cm]{geometry}
\usepackage{graphicx}
\usepackage{booktabs}
\usepackage{float}
\usepackage{caption}
\usepackage{csquotes}
\usepackage[style=numeric-comp,sorting=none,backend=biber,maxbibnames=4]{biblatex}
\usepackage[hidelinks]{hyperref}
\usepackage{cleveref}
\addbibresource{../../latex/references.bib}
\graphicspath{{../../../results/morvan/figures/}{../../../results/figures/}}
\captionsetup{font=small,labelfont=bf}

\title{L band (NISAR) vs C band (Sentinel-1) for clear-cut detection\\
in the Morvan, following the Irstea guide method\\[0.4em]
\large Replication of the Vosges feasibility test -- SERTIT / ICube internship}
\author{Mamadou Galy Diallo}
\date{November 2026}

\begin{document}
\maketitle

\begin{abstract}
\noindent
%%ABSTRACT%%
\end{abstract}

\tableofcontents
\newpage

%=====================================================================
\section{Context and objective}

The SERTIT / ICube internship explores L and P bands (NISAR, BIOMASS) for the detection of deforestation and
landslides. A first test, in the Vosges, compared NISAR (L band, level-2 products generated by ISCE3~\cite{isce3})
and Sentinel-1 (C band~\cite{torres2012sentinel1}) on a 2.4~ha clear-cut. L band separated the cut from the forest
almost perfectly in interferometric coherence (AUC 0.995 vs 0.72 in C band), but its advantage in backscatter
depended on the acquisition geometry. That test had two limits: \textbf{a single cut} and \textbf{a single
Sentinel-1 orbit}.

This replication removes both limits and relies on a French operational method: the Irstea (UMR TETIS)
methodological guide for clear-cut detection~\cite{ose2015guide}, written for the French Ministry of Agriculture on
sylvo-ecoregion G23 \enquote{Morvan et Autunois}. The guide is used both:
\begin{enumerate}
\item as an \textbf{optical reference}: its method is applied to Sentinel-2 over the whole of G23, then checked by
  photo-interpretation;
\item as an \textbf{evaluation standard}: its thresholding method is transposed to radar change images and
  evaluated with its area confusion matrix and its targets ($\leq 10$~\% omission and commission).
\end{enumerate}
L band will be provided in Europe by the Copernicus ROSE-L mission~\cite{esa2020rosel}; the processing chain
developed here transfers directly to it.

\paragraph{Question.} \emph{On several clear-cuts of the Morvan (June--November 2026), and at equal geometry in
both look directions, does NISAR L band separate the cuts from intact forest better than Sentinel-1 C band? And can
an operational method such as the guide's be transposed to radar?}

\paragraph{Hypotheses}, as in the Vosges:
\begin{description}
\item[H1] Coherence of intact forest is higher in L band than in C band.
\item[H2] The coherence loss caused by the cut is therefore better detected in L band.
\item[H3] The backscatter drop after the cut is stronger in L band.
\item[H4] Cross-polarisation is the most sensitive backscatter indicator.
\end{description}

%=====================================================================
\section{Study area and data}

\subsection{Sylvo-ecoregion G23 \enquote{Morvan et Autunois}}

The study area is that of the Irstea methodological guide~\cite{ose2015guide}: sylvo-ecoregion
G23~\cite{ign2011ser}, a low mountain range (about 200 to 900~m) of 3,011~km² across the Nièvre, Yonne,
Côte-d'Or and Saône-et-Loire departments. Forest covers about half of the area (52~\% according to BD Forêt). The
range is known for its conifer plantations, mainly Douglas fir, which are harvested by clear-cutting and are
gradually replacing broadleaved forest (oak, beech). BD Forêt V2~\cite{ign2024bdforet}, read through the
Géoplateforme WFS service, gives 13,494 forest polygons. Keeping only forests with tree cover and poplar stands
(class FF0, \enquote{closed forest without cover}, corresponds to old cuts) gives about 157,800~ha: 89,750~ha of
broadleaves, 49,400~ha of conifers and 18,600~ha of mixed forest.

\subsection{Data}

\begin{table}[htbp]
\centering\small
\caption{Data used in the Morvan (June--November 2026).}
\label{tab-data}
\begin{tabular}{p{2.7cm}p{5.4cm}p{6.7cm}}
\toprule
Sensor & Products & Details \\
\midrule
NISAR (L, 1.26~GHz) & L2 GCOV ($\gamma^0$ HH, HV, RTC) and GUNW (HH coherence at 20~m), produced by
ISCE3~\cite{isce3} & 2 geometries covering the radar window: A173~F25 (ascending) and D66~F65 (descending). Pairs of
12 to 24~days. \\
Sentinel-1 (C, 5.4~GHz) & IW SLC, VV+VH, processed with SNAP~\cite{snap} & Orbits A59 (IW3) and D37 (IW2), 3 bursts
each. 15 dates 12~days apart per orbit, from 01/06 to 18/11/2026, almost all from S1D. \\
Sentinel-2~\cite{drusch2012sentinel2} & L2A, NDVI and SWIR & Reference: guide method over the whole of G23 (tiles
31TEM, 31TEN, 31TFM, 31TFN) and photo-interpretation. \\
IGN & G23 outline, BD Forêt V2 & Forest mask and groups (broadleaves, conifers, mixed). \\
\bottomrule
\end{tabular}
\end{table}

As in the Vosges, download volumes are reduced: NISAR products are read remotely by HDF5 chunks, and only the files
of the useful sub-swath are extracted from the Sentinel-1 archives.

%=====================================================================
\section{Method}

\subsection{The guide method}

The guide~\cite{ose2015guide} detects clear-cuts in six steps: NDVI of two summer images, difference
$\Delta$NDVI, study-area and forest masks, statistical thresholding inside the forest, 1~ha sieve and evaluation by
an area-based confusion matrix. Thresholding uses the mean $m$ and standard deviation $\sigma$ of $\Delta$NDVI in the
forest, with three certainty degrees:
\begin{center}\small
\begin{tabular}{ll}
\toprule
Degree & Interval \\
\midrule
1 (low) & $[m-2\sigma, m-\sigma[$ \\
2 (medium) & $[m-3\sigma, m-2\sigma[$ \\
3 (high) & $< m-3\sigma$ \\
\bottomrule
\end{tabular}
\end{center}
Its specifications are a 1~ha minimum mapping unit and omission and commission errors $\leq 10$~\% by area. The
module \texttt{nisar\_s1.guide\_cr} implements these steps for any change image ($\Delta$NDVI, $\Delta\gamma^0$,
coherence loss), with the direction of change as a parameter.

\subsection{Optical reference}

\paragraph{Guide applied to Sentinel-2.} For each tile, NDVI is the median of the 3 clearest dates of two windows:
\enquote{before} (15/06--08/07/2026) and \enquote{after} (15/09--16/10/2026). The guide thresholds are computed in
the forest of G23.

\paragraph{Qualified reference.} The guide keeps only parcels \enquote{that changed from a wooded state to bare
soil}. The reference is therefore: degree~3, NDVI before $\geq 0.70$, NDVI after $\leq 0.65$, then a 1~ha sieve.

\paragraph{Photo-interpretation.} As recommended by the guide (section 3.9.1), the reference is checked visually:
true colour before and after, and SWIR / NIR / red false colour after, where bare soil and logging residues appear
pink-magenta. Each object receives a verdict: certain cut, probable cut or no bare soil.

\paragraph{\enquote{Forest before} test.} A clear-cut is the loss of a closed stand. The July SWIR reflectance
(B11) inside the object is compared with a 100 to 400~m ring of forest: the object was a closed stand if its B11
does not exceed the ring's 90th percentile.

\subsection{Radar window}

Sentinel-1 processing is limited to 3 bursts of a single sub-swath (RAM). The $20\times20$~km windows were ranked
by the area of reference cuts, then checked in the ASF catalogues: one ascending and one descending Sentinel-1
orbit, each in a single sub-swath, and at least one ascending and one descending NISAR geometry. The selected
window (north-east of G23) allows, unlike the Vosges, a comparison \textbf{at equal geometry in both directions}:
A173 (L) vs A59 (C) ascending, D66 (L) vs D37 (C) descending.

\subsection{Radar processing}

\begin{itemize}
\item \textbf{NISAR}: $\gamma^0$ HH and HV from GCOV (RTC~\cite{small2011flattening}) and HH coherence from GUNW,
  used as delivered by ISCE3 and resampled onto the 20~m Lambert 93 grid.
\item \textbf{Sentinel-1}: same SNAP graphs as in the Vosges. For $\gamma^0$: calibration, Terrain-Flattening and
  geocoding. For coherence: $10\times3$ window, flat-earth and topographic phases removed. The 3 bursts are joined
  by TOPSAR-Deburst, without ESD.
\end{itemize}

\subsection{Change variables}

Acquisitions are classed as \enquote{before} (01/06--12/07), \enquote{after} (10/09--20/10) and \enquote{between}
these two windows, which bracket the optical windows. Cut dates are not known one by one, so the variables are
defined to be independent of the exact date.
\begin{itemize}
\item $\Delta\gamma^0$: mean $\gamma^0$ after $-$ mean $\gamma^0$ before (dB, means in linear power);
\item $\gamma^0$ after (dB);
\item minimum coherence over the \enquote{between} pairs: the pair spanning the cut is one of them, whatever the cut
  date;
\item coherence loss (coherence before $-$ minimum coherence) and coherence after.
\end{itemize}

\subsection{Evaluation}

\paragraph{Classes.} The \enquote{cut} class is made of the forest pixels of the reference cuts. The
\enquote{intact forest} class is BD Forêt forest whose $\Delta$NDVI stays within $m \pm \sigma$, more than 40~m
from any optical detection. Objects discarded by photo-interpretation and a 40~m buffer are excluded from both
classes.

\paragraph{Separability.} For each variable: AUC (Mann-Whitney statistic~\cite{hanley1982auc}, expressed as
$\max(\mathrm{AUC}, 1-\mathrm{AUC})$), Cohen's d~\cite{cohen1988power} and Jeffries-Matusita
distance~\cite{swain1978remote}. 95~\% confidence intervals come from a cluster bootstrap~\cite{kunsch1989bootstrap}:
whole cuts are resampled (not pixels, which are correlated), together with $5\times5$ pixel blocks of forest. With 7
cuts, these intervals are wide.

\paragraph{Checks.} Before the cut, cuts and forest must be indistinguishable. This is checked with the change
between two \enquote{before} dates and with the $\gamma^0$ level before the cut, which reveals a possible
stand-type effect.

\paragraph{Guide transposed to radar.} The $m-k\sigma$ certainty degrees are computed on each radar variable in the
forest of the window, then sieved at 1~ha. The guide's area confusion matrix against the reference is then
computed (omission, commission, kappa), together with the share of cuts found (at least 25~\% of their pixels
detected). The false alarm rate relative to the forest area is also computed: with a reference of a few hectares,
the guide's relative commission is very sensitive to a few hectares of false alarms.

\paragraph{False alarm audit.} The largest radar detections outside any optical detection are photo-interpreted
on a late Sentinel-2 image (November): some may be cuts made after the September optical window.

\paragraph{Dating.} For each reference cut and each geometry, the cross-polarised $\gamma^0$ series is compared
with its pre-July level. The cut is dated between the last normal acquisition and the first acquisition of a
lasting drop of more than 2~dB (two consecutive acquisitions).


%%RESULTS%%

\printbibliography[heading=bibintoc,title={References}]

\appendix
%%APPENDIX%%

\end{document}
